\section{The holomorphic superpotential}
\label{sec:yukawa}

\subsection{Selection rules as a network-flow problem}

Under the unbroken $U(1)^4=S(U(1)^5)$ the low-energy fields carry charges in
$\Z^5/\Z(1,1,1,1,1)$:
\[
\ten_a:\ e_a,\qquad \bfive_{ab}:\ e_a+e_b,\qquad \five_{ab}:\ -(e_a+e_b),\qquad
S_{ab}:\ e_a-e_b,
\]
from $H^1(L_a)$, $H^1(L_a\otimes L_b)$, $H^1(L_a^*\otimes L_b^*)$ and
$H^1(L_a\otimes L_b^*)$ respectively. An operator is neutral when its total
charge $c\in\Z^5$ equals $m(1,\dots,1)$; since singlet charges have entry sum
zero, $m=\tfrac15\sum_ic_i$ is forced, and the operator is forbidden outright if
that is not an integer. Otherwise the singlets must supply the \emph{residual}
$r=m\mathbf1-c$.

\begin{proposition}[Singlet insertions]\label{prop:flow}
Let $E$ be the set of pairs $(a,b)$ for which $S_{ab}$ is present, and $D_E$ the
directed graph on $\{1,\dots,5\}$ with an edge $a\to b$ for each $(a,b)\in E$. A
residual $r$ with $\sum_ir_i=0$ is the charge of a monomial in the available
singlets if and only if $D_E$ carries a nonnegative flow with net outflow $r_a$
at every node $a$. Integer and real feasibility coincide, and the minimal number
of insertions is the value of a transportation problem with costs equal to
shortest-path lengths in $D_E$.
\end{proposition}
\begin{proof}
A monomial $\prod S_{ab}^{n_{ab}}$ has charge $\sum n_{ab}(e_a-e_b)$, the divergence
of the flow $n$. The node--arc incidence matrix is totally
unimodular~\cite{Schrijver1986}, so a feasible real flow implies an integer one.
With unit edge costs and infinite capacities, a min-cost flow decomposes into
paths from sources to sinks, each routed along a shortest path~\cite{AMO1993}.
\end{proof}

Membership is thus decided exactly and with \emph{no bound on the insertion
order}. \st{FORBIDDEN\_BY\_SELECTION\_RULE} is an exact vanishing statement to
all orders in the singlets; an allowed operator is only charge-allowed, and its
value needs the cohomology. For comparison with truncated tables we also
enumerate every monomial up to a given degree, pruning with the min-cost-flow
value as an admissible lower bound, so the enumeration is complete. Two
consequences we use later: a rank bound valid to all orders is the
\emph{structural rank} of the allowed pattern of a Yukawa matrix; and the
singlets that can take a $D$-flat vacuum value at vanishing Fayet--Iliopoulos
terms are exactly those on directed cycles of $D_E$, since $D$-flatness,
$\sum_SQ_a(S)|s|^2=0$, asks for a nonnegative circulation.

\subsection{Reproducing every published operator table}

The standard-model database~\cite{AGLP2012,OxfordCICY} records, for every model,
thirteen operator tables as lists of singlet monomials. We rebuilt all of them
from the charges, for all 2,012 models with complete records (\cref{tab:operators}).

\begin{table}[t]
\centering\small
\begin{tabular}{@{}llr@{}}
\toprule
table & operator & agreement \\
\midrule
\texttt{YukawaUp}, \texttt{YukawaDown}, \texttt{MuTerm} & $\ten\ten\five_H$, $\ten\bfive\bfive_H$, $\five_H\bfive_H$ & 6,036 / 6,036 tables \\
\texttt{HiggsMass} & full Higgs mass matrix & 18,108 / 18,108 cells \\
\texttt{RParVio} & $\bfive^p\five_H$ & 6,036 / 6,036 cells \\
\texttt{Dim4Proton} & $\ten\bfive\bfive$ & 36,216 / 36,216 classes \\
\texttt{Dim5Proton} & $\ten\ten\ten\bfive$ & 60,360 / 60,360 classes \\
\texttt{SuperPotential} & pure-singlet superpotential & 2,012 / 2,012 tables \\
\texttt{DTerms}, \texttt{FITerms} & $D$-terms; FI terms $=\mu_{k_a}(t)$ & 2,012 / 2,012 each \\
\texttt{GMTerm} & $H_uH_d\,f(S,S^\dagger)$ & 2,012 / 2,012 tables \\
\texttt{NeutrinoKahler} & $L^pH_d^\dagger f(S,S^\dagger)$ & 6,036 / 6,036 cells \\
\texttt{Kin5} & $\bfive^{p\dagger}\bfive^q f(S,S^\dagger)$ & 18,108 / 18,108 cells \\
\bottomrule
\end{tabular}
\caption{Every operator table of the published standard-model database,
reproduced from the $U(1)^4$ charges alone for all 2,012 complete models. The
FI-term agreement also checks our intersection numbers on all 13 geometries.
The one flag-level difference in the Yukawa tables is a coupling first allowed
at insertion order six, above the published truncation, where both lists are
empty.}
\label{tab:operators}
\end{table}

Several conventions are not stated in the sources. Each was found from a failing
comparison, diagnosed, then held fixed for the whole database, and the raw
(convention-free) agreement is recorded alongside:
superpotential lists stop at total operator degree eight capped at singlet
order five, Kähler lists at singlet degree three; the Yukawa matrices use the
physical Higgs pair \texttt{PhysHiggs}; families are ordered by bundle index,
which we confirmed by decoding the per-model \texttt{EDownSpec} field; the bare
bilinear is omitted from \texttt{RParVio}; \texttt{Kin5} is indexed
$\bfive^{p\dagger}\bfive^q$; Kähler monomials never contain a trivial
$|S_{ab}|^2$ factor (with such factors only 94 of 2,012 GM tables agree);
\texttt{NeutrinoKahler} lists only monomials containing a holomorphic singlet,
the paper's right-handed neutrino. Two further conventions are more
interesting because they show the tables are not always functions of the charges
alone. \texttt{Dim5Proton} lists an inconsistent subset of the permutations of
its three $\ten$ indices---for one index shape sometimes one ordering of three,
sometimes two---so only permutation classes are meaningful; per class, every
listed cell equals ours. And \texttt{Dim4Proton} omits every entry that involves a
matter $\bfive$ from the pair bundle of the physical down Higgs (charge-identical
to a Yukawa coupling): all 19,450 allowed cells without such a $\bfive$ are
listed and all 164 with one are not.

\paragraph{An internal inconsistency.} Because these are charge tables, two
families from the same line bundle must give identical rows. Using only the
published data---family order from each model's own \texttt{EDownSpec}---all
2,012 published up-type tables pass this test, but \textbf{954 of 2,007
down-type tables fail it}: they contain a zero row where an identical row is
nonzero. Transposing or exchanging labels does not repair it (1,028 failures
either way). Entry by entry on $\Xa$ we reproduce 198/198 up-type and 184/198
down-type cells; the 14 remaining cells lie in exactly five of these
inconsistent tables. \Cref{sec:triple} shows that where the cohomology can
decide, it contradicts the published presentation.

\subsection{Leading-order couplings as triple products}
\label{sec:triple}

At the Abelian locus a Yukawa coupling is the cup product
$H^1(X,L_1)\times H^1(X,L_2)\times H^1(X,L_3)\to H^3(X,\Oc_X)\cong\C$ with
$L_1L_2L_3=\Oc_X$~\cite{AGGHL2010,BBCL2016}. For a hypersurface $X=\{f=0\}$ of
multidegree $q$ in $A=(\PP^1)^4$, the Koszul sequence
$0\to\Oc_A(k-q)\xrightarrow{f}\Oc_A(k)\to\Oc_X(k)\to0$ splits $H^1(X,\Oc(k))$ into
\emph{level 0} (the image of $H^1(A,\Oc(k))$) and \emph{level 1} (mapping onto
$\ker\bigl(H^2(A,\Oc(k-q))\xrightarrow{f}H^2(A,\Oc(k))\bigr)$).

\begin{lemma}[Level-zero vanishing]\label{lem:level0}
If all three classes are at level 0 the coupling vanishes.
\end{lemma}
\begin{proof}
Restriction commutes with cup product, so the product is the restriction of a
class in $H^3(A,\Oc_A)$, which is zero by Künneth since $H^{>0}(\PP^1,\Oc)=0$.
\end{proof}

\begin{proposition}[Level one]\label{prop:level1}
If exactly one class is at level 1, with image $\beta\in H^2(A,\Oc(k_3-q))$, and
the others are restrictions of $\alpha_i\in H^1(A,\Oc(k_i))$, the coupling equals
the ambient product $\alpha_1\alpha_2\beta\in H^4(A,\Oc(-q))\cong\C$.
\end{proposition}
\begin{proof}
Since $H^3(A,\Oc_A)=H^4(A,\Oc_A)=0$ the connecting map
$\delta:H^3(X,\Oc_X)\to H^4(A,\Oc(-q))$ is an isomorphism, and $\delta$ is linear
over restrictions of ambient classes:
$\delta(\alpha|_X\cdot\gamma)=\alpha\cdot\delta\gamma$.
\end{proof}

When two or more classes are at level 1 the connecting map is not
multiplicative, the naive product vanishes only for degree reasons, and the
true coupling needs a lift through $f$~\cite{BBL2017}; such couplings are
returned as \st{UNKNOWN}. (An early version of our code treated them as zero; an
audit showed no recorded coupling was of this type, and a regression test now
pins a genuine level-sum-two triple to \st{UNKNOWN}.) The ambient product is
evaluated on \v{C}ech monomials: on each $\PP^1$ the three exponents must sum to
$(-1,-1)$, the generator of $H^1(\PP^1,\Oc(-2))$.

\paragraph{Complex-structure independence.} If the relevant Koszul map is zero
because its other end is a zero ambient group, the identification of $H^1(X,L)$
with the ambient group holds for \emph{every} smooth member. This was the case
for every one of the 774 bundle slots ever evaluated, so every leading-order
matrix reported is constant over the invariant family: there are no rank-drop
loci at leading order. Couplings where a Koszul map is nonzero are refused.

\paragraph{Downstairs restriction and a convention check.} Each slot is
restricted to the isotypic component of its piece's Wilson-line character minus
the equivariant structure of its bundle, the convention that reproduces 1,948 of
1,948 published spectra. A blocking check requires every slot's isotypic
dimension to equal the number of families the spectrum assigns to that bundle.
The trilinear form must be $\Gam$-covariant with the character of
$H^4(A,\Oc(-q))$, $\prod_i\det(B_i)^{-1}$, which is trivial exactly when the action
descends to the Calabi--Yau quotient. We verified this element by element and
index triple by index triple: \textbf{zero violations}. That single test
validates the Serre-duality bases, the exponent conventions, the Kronecker
ordering and the group lift.

\paragraph{Against the published models.} The published tables record only
which entries are allowed, so we compare our exact leading-order ranks with the
published \emph{structural} rank. For both $Y_u$ and $Y_d$ the ranks agree on
20 of 22 $\Xa$ models. The two $Y_u$ differences are \cref{lem:level0} at work:
charge-allowed couplings whose three classes all sit at level 0, identically zero
though no charge argument can see it. The two $Y_d$ differences are internally
inconsistent published tables; in one, the cohomology gives a $2\times3$ block of
rank two where the published table shows a zero row, and a rank-two $2\times3$
matrix has a zero row in no basis. Finally, \cite[Sec.~6.2]{AGLP2012} proves that
without equivariant blocks $Y_u$ cannot have rank one; our computed $Y_u$ ranks
over the 22 models are 0 or 2 and never 1, recovering that theorem from the
cohomology rather than from the group theory.

\subsection{The Yukawa sector of the downstairs models}

We now apply the engine to all 21 orbits of \cref{sec:downstairs}---the 11
additional ones and, as a baseline, the 10 published ones---through the
identical pipeline, over every accepted configuration (192--336 per orbit) and
every choice of physical Higgs pair. Singlet spectra are derived from first
principles; validated on the published $\Xa$ models they agree exactly for 21
of 22, and the 22nd differs only on two singlets whose published multiplicity is
an unevaluated rank expression that our equivariant rank computation evaluates
(to zero).

Besides the Yukawa matrices we evaluate every dangerous operator that is a
genuine triple product: dimension-four $R$-parity violation at singlet order
zero (its $QLd^c$, $LLe^c$ and $u^cd^cd^c$ pieces), and the one-singlet terms
$SH_uH_d$ and $SLH_u$. Higher-order terms are decided by the selection rule only.
We then classify with logical categories, applied identically to both groups:
\st{STRONG\_CANDIDATE} requires a nonzero top Yukawa at the Abelian locus,
rank three reachable for $Y_u$ and $Y_d$ at some singlet order, and
dimension-five proton decay, dimension-four and bilinear $R$-parity violation
forbidden to all orders; \st{PHENOMENOLOGICALLY\_PROBLEMATIC} requires an
unconditional exact obstruction---no up- or no down-type Yukawa at any
perturbative order, or dimension-four $R$-parity violation present at the
Abelian locus; everything else is \st{UNKNOWN}. An orbit is strong if some
choice is, problematic if every choice is.

\begin{table}[t]
\centering\small
\setlength{\tabcolsep}{4pt}
\begin{tabular}{@{}llrcccccl@{}}
\toprule
orbit & published & choices & top at LO & cap $Y_u$ & cap $Y_d$ & QQQL forbidden & RPV$_4$ at LO & verdict \\
\midrule
A01 & -- & 192 & -- & 0 & 3 & all & -- & problematic \\
A02 & -- & 720 & -- & 0 & 1, 3 & all & 288 & problematic \\
A03 & -- & 192 & -- & 0 & 1 & all & -- & problematic \\
A04 & -- & 720 & -- & 3 & 3 & none & 288 & \textbf{unknown} \\
A05 & -- & 192 & -- & 0 & 0 & all & -- & problematic \\
A06 & -- & 720 & -- & 3 & 3 & none & -- & \textbf{unknown} \\
A07 & -- & 720 & -- & 0 & 1 & all & 288 & problematic \\
A08 & -- & 192 & -- & 0 & 0 & all & -- & problematic \\
A09 & -- & 192 & -- & 0 & 0 & all & -- & problematic \\
A10 & -- & 720 & -- & 0 & 0 & all & -- & problematic \\
A11 & -- & 720 & -- & 0 & 0 & all & -- & problematic \\
\midrule
P01 & 1 & 192 & -- & 0 & 0 & all & -- & problematic \\
P02 & 2 & 192 & all & 2 & 0 & all & -- & problematic \\
P03 & 3--6 & 384 & -- & 0 & 0, 1 & all & -- & problematic \\
P04 & 7 & 192 & -- & 0 & 0 & all & -- & problematic \\
P05 & 8 & 192 & -- & 0 & 0 & all & -- & problematic \\
P06 & 9--12 & 384 & -- & 0, 2 & 2, 3 & none & 192 & problematic \\
P07 & 13 & 192 & -- & 0 & 1 & none & -- & problematic \\
P08 & 14 & 192 & -- & 3 & 3 & none & -- & unknown \\
P09 & 15--18 & 384 & -- & 0, 3 & 0, 1, 3 & 288 & 96 & unknown \\
P10 & 19--22 & 384 & 192 & 1, 2 & 1, 2 & none & 96 & unknown \\
\bottomrule
\end{tabular}
\caption{The downstairs models on $\Xa/(\Z_2\times\Z_2)$. ``Choices'' counts
configurations times physical-Higgs choices; ``cap'' is the largest rank the
selection rule allows at any singlet order (sets list the values over choices);
counts are numbers of choices. ``Published'' gives the published model numbers
in each orbit.}
\label{tab:ranking}
\end{table}

\Cref{tab:ranking} gives the result, and it answers the question posed for the
new models. \textbf{No orbit on this quotient, published or new, is a strong
candidate.} Nine of the eleven new orbits are problematic and two unknown; seven
of the ten published orbits are problematic and three unknown. The new orbits
fail for the same reasons as the published ones, at a comparable rate. In
detail, for the eleven new orbits:
\begin{itemize}[leftmargin=*]
\item none has a nonzero top Yukawa at the Abelian locus (two published orbits
  do); in nine the up-type Yukawa is forbidden to \emph{all} perturbative orders;
\item A04 and A06 are the only new orbits whose up-type Yukawa can reach rank
  three, and for both the rank-three pattern survives restricting to singlets
  that can take $D$-flat values;
\item dimension-five proton decay is forbidden to all orders in nine of eleven,
  in every configuration;
\item dimension-four $R$-parity violation is computed nonzero at the Abelian
  locus in 288 of 720 Higgs choices of A02, A04 and A07, and the one-singlet
  bilinear coupling $SLH_u$ is nonzero in every choice of A04 and A07;
\item in A08, $SLH_u$ is charge-allowed but \emph{vanishes cohomologically} in
  every configuration: an allowed-but-vanishing dangerous operator.
\end{itemize}

The \st{PROBLEMATIC} verdicts are exact but perturbative: masses forbidden to all
orders by the $U(1)^4$ symmetries can still be generated non-perturbatively where
those symmetries are anomalous. The \st{UNKNOWN} verdicts rest on the
coefficients of higher-order couplings, on the singlet vacuum, and on the
normalisation of the physical couplings, none of which is computed here.
